% !TeX root = ../main.tex
\section{Zusammenfassung}

Elektrische Organentladungen (EODs) bieten eine nicht-invasive Möglichkeit, das Verhalten von Zitteraalen über Zeitskalen hinweg zu untersuchen, die durch direkte Beobachtung nur schwer zu erfassen sind.
Diese Arbeit untersucht Langzeitmuster in den EOD-Wellenformen und der Aktivität von in Gefangenschaft gehaltenen \textit{Electrophorus voltai} anhand kontinuierlicher Mehr-Elektroden-Aufzeichnungen aus dem Museum für Naturkunde Berlin.
Der Datensatz umfasst mehr als 13 Millionen akzeptierte EODs, die zwischen 2023 und 2026 aufgezeichnet wurden.
Die Wellenformen wurden mithilfe einer Kombination aus robuster Hauptkomponentenanalyse und Random-Forest-Klassifizierung klassifiziert, und die Impulsaktivität wurde über mehrere Zeitskalen hinweg sowie in Bezug auf die Fütterung und die Position im Becken analysiert.
Es wurden drei reproduzierbare Wellenformklassen identifiziert: normale, breite und Doppelimpulse.
Die Gesamtimpulsaktivität zeigte eine deutliche fütterungsbezogene Reaktion und stieg von etwa 1,34~Hz außerhalb der Fütterungsfenster auf 5,03~Hz während der Fütterung an.
Die Aktivität der breiten Impulse zeigte eine ähnliche Reaktion, während die Doppelimpulsaktivität niedrig blieb und keinen vergleichbaren Anstieg der mittleren Frequenz aufwies.
Über die gesamte Langzeitaufzeichnung hinweg trat der stärkste Anstieg der Spezialimpulsaktivität während eines bestimmten Zeitraums Ende 2025 auf und nicht als wiederkehrendes Muster im Kalendermonat.
Dieser Anstieg fiel mit einer Verschiebung der grobskaligen Raumnutzung in Richtung des dunkleren Bereichs des Beckens sowie mit ungewöhnlichen Beobachtungen einer erhöhten Nähe zwischen den Tieren zusammen.
Insgesamt stehen diese Veränderungen im Einklang mit einer Veränderung des Sozialzustands und lassen die Möglichkeit eines Fortpflanzungskontextes vermuten; die vorliegenden Daten belegen jedoch weder ein Fortpflanzungsereignis noch lassen sie sich bestimmte Impulse als Fortpflanzungssignale identifizieren.
Die Ergebnisse bilden eine langfristige elektrophysiologische Basislinie für in Gefangenschaft lebende \textit{E. voltai} und bieten einen analytischen Rahmen für zukünftige Vergleiche mit Aufzeichnungen aus der natürlichen Umgebung.

\newpage